\chapter{Stato dell'arte}
\label{ch:stato-dell-arte}

Uno degli studi pionieristici della \textit{sentiment analysis} è quello di \textcite{tetlock2007giving}.
In particolare, Tetlock ha quantificato il ``pessimismo'' degli articoli del \textit{Wall Street Journal} e ha rilevato che livelli elevati di tono negativo nella stampa anticipano pressioni al ribasso sui prezzi azionari, seguite da un ritorno ai fondamentali di mercato.
Questo risultato ha fornito evidenza empirica che il sentiment estratto dai testi---ad esempio il tono pessimista o ottimista delle notizie---può influenzare il comportamento di mercato.

\section{Social media e sentiment degli investitori}

Con la diffusione dei social media e delle piattaforme online negli anni 2010, le fonti di dati sul sentiment si sono ampliate.
Gli investitori, in particolare quelli retail, hanno iniziato a esprimere opinioni su forum finanziari, Twitter e blog, generando un flusso di testo ricco di informazioni sul sentiment di mercato in tempo reale.

La letteratura di questo periodo ha analizzato in che misura il sentiment influenzi i rendimenti azionari, spesso in tensione con la teoria finanziaria tradizionale, che tende a escludere fattori emotivi.
\textcite{mcgurk2020stock} illustrano chiaramente questo cambio di paradigma: analizzando i post sui social media, gli autori costruiscono una misura quantitativa del sentiment degli investitori e rilevano un effetto positivo e significativo sui rendimenti anomali delle azioni.

Un caso emblematico che ha catalizzato l'interesse sull'analisi del sentiment è la pandemia di COVID-19.
\textcite{duan2021covid}, tra gli altri, sviluppano indici di sentiment legati al COVID-19 utilizzando milioni di testi provenienti sia dai media ufficiali sia dai social in Cina.

\section{L'avvento dei modelli Transformer e di FinBERT}

Alla fine degli anni 2010, l'introduzione di BERT \parencite{devlin2018bert} ha rivoluzionato l'NLP anche in ambito finanziario, grazie a rappresentazioni linguistiche contestuali di alta qualità.
Tuttavia, l'adattamento al linguaggio tecnico della finanza ha richiesto modelli specializzati.
FinBERT, proposto da \textcite{araci2019finbert}, è stato tra i primi modelli pre-addestrati su testi finanziari e successivamente fine-tunati per la sentiment analysis, superando di circa 15 punti percentuali i metodi precedenti su dataset come \textit{Financial PhraseBank}.

Versioni successive, come quella di \textcite{huang2022finbert}, addestrate su ampi corpora finanziari, hanno raggiunto fino all'89,7\% di accuratezza nel rilevare frasi negative, migliorando la comprensione delle sfumature linguistiche rispetto ai modelli generici.

FinBERT è oggi uno standard per task NLP finanziari: analisi del tono nei report, classificazione di notizie, estrazione di informazioni da bilanci.
\textcite{gu2024finbert} mostrano che un modello FinBERT+LSTM migliora la previsione dei trend azionari rispetto a LSTM tradizionali, evidenziando il valore della specializzazione di dominio.

Infine, l'emergere di benchmark specifici (es.\ FiQA, Financial PhraseBank) ha supportato valutazioni comparative, mentre estensioni recenti si concentrano su analisi a livello di entità \parencite{tang-etal-2023-finentity}.

\section{Modelli di linguaggio generali vs.\ specifici in finanza}

Nel 2022--2023, l'attenzione si è spostata sui LLM generici come ChatGPT (GPT-3.5/4), capaci di interpretare testi finanziari nonostante l'assenza di addestramento settoriale.
\textcite{lopezlira2024chatgptforecaststockprice} mostrano che ChatGPT può prevedere la direzione dei prezzi azionari basandosi sui titoli delle notizie.
Anche in modalità zero-shot, modelli come GPT-3.5/4 riescono a svolgere sentiment analysis su documenti finanziari, pur con limiti su contesti specialistici o eventi recenti \parencite{bdcc8110143}.

Questi risultati hanno spinto verso LLM specializzati.
BloombergGPT \parencite{wu2023bloomberggptlargelanguagemodel}, con 50 miliardi di parametri e un corpus di 363 miliardi di token finanziari, è tra i primi LLM su larga scala dedicati alla finanza, con maggiore precisione su prompt settoriali grazie alla conoscenza approfondita del dominio.

In parallelo, il progetto open-source FinGPT \parencite{yang2023fingptopensourcefinanciallarge} punta a rendere i LLM finanziari accessibili, attraverso fine-tuning di modelli open-source su dati finanziari pubblici.
Con un approccio \textit{data-centric}, FinGPT abilita analisi di sentiment, generazione automatica di report e sviluppo low-code, pur operando con risorse più contenute rispetto ai modelli proprietari.

\section{Instruction tuning e preference optimization}

Una direzione che si è consolidata dal 2023 è l'\textit{instruction tuning} di LLM open-weight su compiti finanziari.
Il progetto PIXIU \parencite{xie2023pixiu} introduce FinMA, un modello basato su LLaMA e affinato su 136.000 istruzioni finanziarie multi-task, insieme al primo benchmark standardizzato di valutazione (FLARE).
I modelli instruction-tuned di dimensioni contenute eccellono nei task di sentiment classification (Financial PhraseBank, FiQA-SA), pur generalizzando peggio dei grandi modelli proprietari su compiti di ragionamento complesso.

Il paradigma più recente supera il fine-tuning supervisionato (SFT) con tecniche di \textit{preference alignment}.
FinDPO \parencite{iacovides2025findpo} applica la Direct Preference Optimization a un LLM causale per la sentiment analysis finanziaria, superando in media dell'11\% i modelli SFT esistenti sui benchmark standard.
Di particolare interesse per questo progetto è la conversione \textit{logit-to-score} proposta dagli autori: le predizioni discrete di sentiment vengono trasformate in punteggi continui e ordinabili (probabilità), direttamente integrabili in strategie di portafoglio---un'idea affine al \textit{sentiment score} continuo adottato nella presente proposta.
Nelle simulazioni riportate, la strategia basata su FinDPO mantiene rendimenti positivi anche in presenza di costi di transazione realistici.

Più in generale, diversi lavori del 2024--2025 mostrano che LLM open-weight recenti (LLaMA~3, Qwen, DeepSeek), anche con fine-tuning leggero, superano stabilmente FinBERT nella classificazione del sentiment finanziario, rendendo lo scoring di qualità accessibile senza infrastrutture proprietarie.

\section{Benchmark e valutazione sistematica}

La proliferazione di modelli ha reso centrale il problema della valutazione.
FinBen \parencite{xie2024finben} è il primo benchmark open-source su larga scala per LLM finanziari: 36 dataset su 24 task, che coprono estrazione di informazioni, analisi testuale, question answering, generazione, gestione del rischio, previsione e decision-making, incluse valutazioni di agenti e di pipeline RAG.
I risultati confermano che i LLM eccellono nell'analisi testuale (incluso il sentiment) ma faticano su ragionamento avanzato e previsione.

In parallelo emergono cautele metodologiche.
\textcite{bdcc8110143} mostrano che un semplice modello di regressione logistica ha superato FinBERT e GPT-4 in accuratezza su un task di sentiment, segnalando che maggiore complessità non implica necessariamente migliori performance.
Inoltre, benchmark storici come Financial PhraseBank risultano ormai prossimi alla saturazione e potenzialmente presenti nei dati di pre-training dei modelli più recenti, sollevando dubbi di contaminazione che spingono verso dataset di valutazione nuovi, a livello di entità \parencite{tang-etal-2023-finentity} e temporalmente successivi al cutoff dei modelli.

\section{Agenti LLM e validità empirica delle strategie}

Dal 2024 la frontiera si è spostata dagli score di sentiment ai sistemi \textit{agentici}, in cui più LLM con ruoli specializzati collaborano al processo decisionale.
TradingAgents \parencite{xiao2025tradingagents} simula una società di trading con agenti analisti (fondamentale, sentiment, tecnico), ricercatori bull/bear in dibattito, trader e risk manager, riportando miglioramenti su rendimento cumulato, Sharpe ratio e drawdown massimo.
Framework analoghi (FinMem, FinCon, FinRobot, MarketSenseAI) integrano memoria episodica, retrieval di notizie e analisi multi-fonte.

A questi risultati promettenti si contrappone una crescente letteratura critica sulla loro validità empirica.
\textcite{li2025finsaber} propongono FINSABER, un framework di backtesting che valuta le strategie LLM su oltre vent'anni e più di cento titoli, controllando \textit{survivorship bias}, \textit{look-ahead bias} e \textit{data-snooping}: i vantaggi riportati in letteratura si riducono drasticamente su orizzonti lunghi e universi ampi, con strategie troppo conservative nei mercati rialzisti e troppo aggressive in quelli ribassisti.
Un'insidia specifica dei LLM è inoltre il look-ahead bias ``parametrico'': un modello con cutoff di addestramento recente ha già osservato gli anni coperti dal backtest, e questa conoscenza---incorporata nei pesi---è invisibile ai controlli sulla pipeline dei dati.
Queste evidenze rendono il disegno sperimentale del backtesting (Capitolo~\ref{ch:descrizione-progetto}) un aspetto metodologicamente critico, non un dettaglio implementativo.

\section{Interazioni tra notizie, grafi e diffusione}

Un ulteriore filone, direttamente rilevante per la proposta di diffusione semantica di questo progetto, abbandona l'ipotesi che le notizie agiscano in modo indipendente.
\textcite{wan2021sentimentdiffusion} mostrano che il sentiment si diffonde lungo reti di notizie tra aziende collegate, con effetti misurabili sui movimenti di mercato; \textcite{wang2024finin} introducono FININ, che modella esplicitamente le interazioni semantiche tra notizie per la previsione dei mercati.
Più recentemente, architetture a grafo con meccanismi di diffusione dell'informazione (graph neural network su relazioni tra titoli, reti di \textit{sentiment spillover} stimate via transfer entropy) integrano segnali testuali e struttura relazionale del mercato, confermando che la propagazione cross-asset dell'informazione è un segnale predittivo autonomo.
La formulazione continua a derivate parziali proposta nel Capitolo~\ref{ch:descrizione-progetto} si colloca in questo filone, ma con un approccio interpretabile e ancora poco documentato in letteratura.

\section{Sintesi e tendenze}

Nuove metriche ampliano infine le prospettive oltre la polarità del tono: \textcite{cao2025hypeindexnlpdrivenmeasure} propongono l'Hype Index, che misura l'attenzione mediatica anziché il sentiment.
Analizzando l'S\&P~100, dimostrano che un picco di copertura può anticipare volatilità e movimenti di mercato, sottolineando il valore dell'``intensità'' informativa come fattore predittivo.

In sintesi, dal 2000 al 2026 l'evoluzione è passata dalle analisi lessicali (Tetlock) ai modelli specializzati (FinBERT), ai LLM generici e finanziari (BloombergGPT, FinGPT), fino a instruction tuning e preference optimization (FinMA, FinDPO), sistemi agentici (TradingAgents) e benchmark olistici (FinBen).
Le tendenze attuali puntano su punteggi continui integrabili in strategie, valutazione rigorosa e priva di bias temporali, analisi a livello di entità e modelli che considerano non solo il sentiment, ma anche attenzione mediatica, interazioni tra notizie e struttura relazionale del mercato.
